% 电路理论基础公式汇总 - 修复版
\chapter{电路理论基础公式汇总}
\thispagestyle{plain}  % 恢复页码样式

\section{第一章：电路元件与电路基本定律}

\subsection{电流、电压与电功率}
\begin{enumerate}
    \item 电流定义：$i=\frac{dq}{dt} \formulaNote{（单位：A，库/秒）}$
    \item 电压定义：$u_{ab}=\frac{dW_{ab}}{dq} \formulaNote{（单位：V，焦/库）}$
    \item 电源电动势：$e=\frac{dW}{dq} \formulaNote{（单位：V，非静电力做功）}$
    \item 电功率定义：$p=\frac{dW}{dt}=ui \formulaNote{（关联参考方向，单位：W）}$
    \item 电路能量计算：$W(t)=\int_{t_0}^tp(\xi)d\xi=\int_{t_0}^tu(\xi)i(\xi)d\xi \formulaNote{（单位：J）}$
\end{enumerate}

\subsection{电阻元件}
\begin{itemize}
    \item 欧姆定律：$u=Ri \quad \text{或} \quad i=Gu \formulaNote{（关联参考方向）}$
    \item 电导与电阻关系：$RG=1 \formulaNote{（G为电导，单位：S）}$
    \item 电阻功率：$p=ui=Ri^2=Gu^2 \formulaNote{（耗能元件，$p\ge0$）}$
    \item 电阻总耗能：
    \[
    W=\int_{-\infty}^tp(\xi)d\xi=\int_{-\infty}^tu(\xi)i(\xi)d\xi = R\int_{-\infty}^ti^2(\xi)d\xi=G\int_{-\infty}^tu^2(\xi)d\xi
    \]
\end{itemize}

\subsection{电容元件}
\begin{itemize}
    \item 电容基本公式：$C=\frac{\varepsilon A}{d} \formulaNote{（平行板电容，$\varepsilon$为介电常数）}$
    \item 电荷-电压关系：$q=Cu \formulaNote{（单位：F，库/伏）}$
    \item 电流-电压微分关系：$i=\frac{dq}{dt}=C\frac{du}{dt} \formulaNote{（直流时$i=0$，等效开路）}$
    \item 电荷积分形式：$q(t)=\int_{-\infty}^ti(\xi)d\xi$
    \item 电压积分形式：$u(t)=\frac{1}{C}\int_{-\infty}^ti(\xi)d\xi$
    \item 初始条件下电压：$u(t)=u(t_0)+\frac{1}{C}\int_{t_0}^ti(\xi)d\xi \formulaNote{（$u(t_0)$为初始电压）}$
    \item 电容储能：$W_e(t)=\frac{1}{2}Cu^2=\frac{q^2}{2C} \formulaNote{（单位：J，仅与当前电压有关）}$
\end{itemize}

\subsection{电感元件}
\begin{itemize}
    \item 磁链-电流关系：$\Psi=Li \formulaNote{（L为电感，单位：H，韦/安）}$
    \item 电压-电流微分关系：$u=\frac{d\Psi}{dt}=L\frac{di}{dt} \formulaNote{（直流时$u=0$，等效短路）}$
    \item 电流积分形式：$i(t)=\frac{1}{L}\int_{-\infty}^tu(\xi)d\xi$
    \item 初始条件下电流：$i(t)=i(t_0)+\frac{1}{L}\int_{t_0}^tu(\xi)d\xi \formulaNote{（$i(t_0)$为初始电流）}$
    \item 电感储能：$W_m=\frac{1}{2}Li^2=\frac{\Psi^2}{2L} \formulaNote{（单位：J，仅与当前电流有关）}$
\end{itemize}

\subsection{独立电源}
\begin{itemize}
    \item 电压源端口特性：$u=u_s \formulaNote{（电压与电流无关，置零等效短路）}$
    \item 电流源端口特性：$i=i_s \formulaNote{（电流与电压无关，置零等效开路）}$
\end{itemize}

\subsection{受控电源（线性）}
\begin{itemize}
    \item 电压控制电压源（VCVS）：$\begin{cases}u_s = \alpha u_c \\i_c = 0\end{cases} \formulaNote{（$\alpha$：电压放大系数）}$
    \item 电流控制电压源（CCVS）：$\begin{cases}u_s = r i_c \\u_c = 0\end{cases} \formulaNote{（$r$：转移电阻，$\Omega$）}$
    \item 电压控制电流源（VCCS）：$\begin{cases}i_s = g u_c \\i_c = 0\end{cases} \formulaNote{（$g$：转移电导，S）}$
    \item 电流控制电流源（CCCS）：$\begin{cases}i_s = \beta i_c \\u_c = 0\end{cases} \formulaNote{（$\beta$：电流放大系数）}$
\end{itemize}

\subsection{基尔霍夫定律}
\begin{itemize}
    \item 基尔霍夫电流定律（KCL）：$\sum i_k=0 \formulaNote{（节点电流守恒，流出为正）}$
    \item 广义KCL：$\sum_{k\in S} i_k=0 \formulaNote{（闭合边界电流守恒）}$
    \item 基尔霍夫电压定律（KVL）：$\sum u_k=0 \formulaNote{（回路电压守恒，绕行方向一致为正）}$
    \item KVL推广：$\sum u_{\text{电压降}}=\sum u_{\text{电压升}}$
\end{itemize}

\section{第十章：二端口网络核心公式}

\subsection{导纳参数（Y参数，短路参数）}
\begin{itemize}
    \item 代数方程：
    \[
    \begin{cases} 
    \dot{I}_1 = Y_{11}\dot{U}_1 + Y_{12}\dot{U}_2 \\ 
    \dot{I}_2 = Y_{21}\dot{U}_1 + Y_{22}\dot{U}_2 
    \end{cases}
    \]
    \item 矩阵形式：$\begin{bmatrix} \dot{I}_1 \\ \dot{I}_2 \end{bmatrix} = \begin{bmatrix} Y_{11} & Y_{12} \\ Y_{21} & Y_{22} \end{bmatrix}\begin{bmatrix} \dot{U}_1 \\ \dot{U}_2 \end{bmatrix}$
    \item 参数计算：
    \[
    \begin{aligned}
    Y_{11} &= \left. \frac{\dot{I}_1}{\dot{U}_1} \right|_{\dot{U}_2 = 0}, & Y_{12} &= \left. \frac{\dot{I}_1}{\dot{U}_2} \right|_{\dot{U}_1 = 0} \\
    Y_{21} &= \left. \frac{\dot{I}_2}{\dot{U}_1} \right|_{\dot{U}_2 = 0}, & Y_{22} &= \left. \frac{\dot{I}_2}{\dot{U}_2} \right|_{\dot{U}_1 = 0}
    \end{aligned}
    \]
    \item 互易/对称条件：$Y_{12}=Y_{21}$（互易）；$Y_{12}=Y_{21}$且$Y_{11}=Y_{22}$（对称）
\end{itemize}

\subsection{阻抗参数（Z参数，开路参数）}
\begin{itemize}
    \item 代数方程：
    \[
    \begin{cases} 
    \dot{U}_1 = Z_{11}\dot{I}_1 + Z_{12}\dot{I}_2 \\ 
    \dot{U}_2 = Z_{21}\dot{I}_1 + Z_{22}\dot{I}_2 
    \end{cases}
    \]
    \item 矩阵形式：$\begin{bmatrix} \dot{U}_1 \\ \dot{U}_2 \end{bmatrix} = \begin{bmatrix} Z_{11} & Z_{12} \\ Z_{21} & Z_{22} \end{bmatrix}\begin{bmatrix} \dot{I}_1 \\ \dot{I}_2 \end{bmatrix}$
    \item 参数计算：
    \[
    \begin{aligned}
    Z_{11} &= \left. \frac{\dot{U}_1}{\dot{I}_1} \right|_{\dot{I}_2 = 0}, & Z_{12} &= \left. \frac{\dot{U}_1}{\dot{I}_2} \right|_{\dot{I}_1 = 0} \\
    Z_{21} &= \left. \frac{\dot{U}_2}{\dot{I}_1} \right|_{\dot{I}_2 = 0}, & Z_{22} &= \left. \frac{\dot{U}_2}{\dot{I}_2} \right|_{\dot{I}_1 = 0}
    \end{aligned}
    \]
    \item 互易/对称条件：$Z_{12}=Z_{21}$（互易）；$Z_{12}=Z_{21}$且$Z_{11}=Z_{22}$（对称）
\end{itemize}

\subsection{传输参数（A参数，T参数）}
\begin{itemize}
    \item 代数方程：
    \[
    \begin{cases} 
    \dot{U}_1 = A_{11}\dot{U}_2 + A_{12}(-\dot{I}_2) \\ 
    \dot{I}_1 = A_{21}\dot{U}_2 + A_{22}(-\dot{I}_2) 
    \end{cases}
    \]
    \item 矩阵形式：$\begin{bmatrix} \dot{U}_1 \\ \dot{I}_1 \end{bmatrix} = \begin{bmatrix} A_{11} & A_{12} \\ A_{21} & A_{22} \end{bmatrix}\begin{bmatrix} \dot{U}_2 \\ -\dot{I}_2 \end{bmatrix}$
    \item 互易/对称条件：$A_{11}A_{22}-A_{12}A_{21}=1$（互易）；$A_{11}A_{22}-A_{12}A_{21}=1$且$A_{11}=A_{22}$（对称）
    \item Z转A参数：$A_{11}=\frac{Z_{11}}{Z_{21}},\ A_{12}=\frac{Z_{11}Z_{22}-Z_{12}Z_{21}}{Z_{21}},\ A_{21}=\frac{1}{Z_{21}},\ A_{22}=\frac{Z_{22}}{Z_{21}}$
\end{itemize}

\subsection{混合参数（H参数）}
\begin{itemize}
    \item 代数方程：
    \[
    \begin{cases} 
    \dot{U}_1 = H_{11}\dot{I}_1 + H_{12}\dot{U}_2 \\ 
    \dot{I}_2 = H_{21}\dot{I}_1 + H_{22}\dot{U}_2 
    \end{cases}
    \]
    \item 矩阵形式：$\begin{bmatrix} \dot{U}_1 \\ \dot{I}_2 \end{bmatrix} = \begin{bmatrix} H_{11} & H_{12} \\ H_{21} & H_{22} \end{bmatrix}\begin{bmatrix} \dot{I}_1 \\ \dot{U}_2 \end{bmatrix}$
    \item 互易/对称条件：$H_{12}=-H_{21}$（互易）；$H_{12}=-H_{21}$且$H_{11}H_{22}-H_{12}H_{21}=1$（对称）
\end{itemize}

\subsection{二端口等效电路}
\begin{itemize}
    \item Z参数T形等效：$Z_1=Z_{11}-Z_{12},\ Z_2=Z_{22}-Z_{21},\ Z_3=Z_{12}$
    \item Y参数Π形等效：$Y_1=Y_{11}+Y_{12},\ Y_3=-Y_{21}=-Y_{12},\ Y_2=Y_{22}+Y_{21}$
    \item 级联总传输矩阵：$\mathbf{A}_{\text{总}} = \mathbf{A}_1 \cdot \mathbf{A}_2 \cdot \dots \cdot \mathbf{A}_n$
\end{itemize}

% 补充通用说明，提升实用性
\section{公式使用说明}
\begin{enumerate}
    \item 所有公式默认采用\textbf{关联参考方向}（电压正极流入，电流流出），非关联方向需加负号；
    \item 交流量用$\dot{U}$、$\dot{I}$表示相量，直流量用$U$、$I$表示；
    \item 互易条件仅适用于无受控源的二端口网络，对称条件需电气结构对称；
    \item 电容/电感的积分公式中，$t=-\infty$表示初始状态为零（未充电/无电流）。
\end{enumerate}